\documentclass[11pt,a4paper]{article}
% Estilo comum dos PDFs individuais da Sprint 1.
\usepackage[utf8]{inputenc}
\usepackage[T1]{fontenc}
\usepackage[brazilian]{babel}
\usepackage[a4paper,margin=2.35cm]{geometry}
\usepackage{xcolor}
\usepackage{hyperref}
\usepackage{enumitem}
\usepackage{listings}
\usepackage{parskip}
\usepackage{needspace}
\definecolor{azul}{RGB}{20,63,112}
\hypersetup{colorlinks=true,linkcolor=azul,urlcolor=azul}
\setlist[itemize]{leftmargin=*,itemsep=3pt}
\lstset{basicstyle=\ttfamily\small,frame=single,breaklines=true,columns=fullflexible}
\sloppy
\newcommand{\codigo}[1]{\texttt{#1}}
\newcommand{\titulo}[2]{\begin{center}{\LARGE\bfseries #1\par}\vspace{4pt}{\large #2\par}\vspace{8pt}\rule{0.8\textwidth}{0.4pt}\par\vspace{5pt}{\small Matheus Moreli Durão Silva -- Matrícula 12621CBS217\par Cibersegurança - Bacharelado - Noturno\par}\end{center}}
\newcommand{\cardfonte}[1]{\par\smallskip\noindent\textit{Referência: }\href{#1}{Trello - enunciado e checklist do card}.}

\begin{document}\titulo{S1\_16 - Git e GitHub}{Registrar a evolução do trabalho}
\section*{Entregável correto}
O card pede um repositório no GitHub com o projeto e seu histórico de commits. Git não substitui o projeto: ele registra versões dele. GitHub é o serviço remoto onde esse repositório pode ser publicado e compartilhado.
\section*{Ciclo básico de versionamento}
Um repositório local nasce com \codigo{git init}. \codigo{git status} mostra o que mudou. \codigo{git add arquivo} seleciona mudanças para a próxima versão; \codigo{git commit -m "mensagem"} cria o marco histórico. Para publicar, configura-se um remoto e usa-se \codigo{git push}. \codigo{git pull} traz mudanças que já existem no remoto. \codigo{git log} permite ler a sequência de commits.
\begin{lstlisting}
git status
git add app.py dados/inventario.json
git commit -m "feat: cadastrar ativos"
git push origin main
\end{lstlisting}
Uma mensagem de commit deve dizer o que mudou: ``feat: cadastrar ativos'' é melhor que ``alteracoes''. Cada commit pequeno permite entender, revisar e recuperar uma etapa sem misturar assuntos.
\section*{Evidência deste projeto}
O repositório publicado é \href{https://github.com/matheus-moreli/inventario-ativos-vulnerabilidades}{inventario-ativos-vulnerabilidades}. O histórico contém commits de funcionalidades e de documentação. Clonar um repositório com \codigo{git clone URL} cria uma cópia local; abrir \codigo{git log} nela mostra as versões anteriores.
\section*{Checklist demonstrado}
Conceito de controle de versão, instalação/configuração, inicialização, commits e mensagens, \codigo{push/pull} e histórico. A conferência prática é abrir o repositório, localizar os arquivos do projeto e verificar que há mais de um commit com propósito identificável.
\cardfonte{https://trello.com/c/vUu4H8Bq}
\end{document}
